%------------------------------------------------------------------------------%
\section{Definição}

Séries são somas com um número infinito de termos. As primeiras questões que
precisamos responder é o que exatamente isso significa e quando essa soma faz
sentido. Para isso vamos utilizar as ideias construídas no estudo das
sequências numéricas.

%------------------------------------------------------------------------------%
\begin{definicao}{Série Numérica}{def:serie-numerica}
  \emph{Série Numérica}\index{Série!numérica}\index{Série!termo geral}
  é a soma dos infinitos termos de uma sequência numérica, $(a_n)$,
  \[
  	\sum_{n=1}^\infty a_n = a_1 + a_2 + a_3 + \cdots
 	\]
  Cada termo $a_n$ é chamado de \emph{Termo Geral}\index{Série!termo geral} da série.
\end{definicao}
%------------------------------------------------------------------------------%

Para determinarmos se essa soma faz sentido precisamos construir um critério de convergência,
para isso, vamos construir outra sequência numérica associada a série.

%------------------------------------------------------------------------------%
\begin{definicao}{Sequência de Somas Parciais}{def:somas-parciais}
  Seja $(a_n)$ uma sequência de números reais, sua
  Sequência de Somas Parciais\index{Somas parciais}\index{Sequência!somas parciais}
  é definida como
  \[
    S_n = a_1 + a_2 + \cdots + a_n = \sum_{k=1}^n a_k   \,.
  \]
\end{definicao}
%------------------------------------------------------------------------------%

Podemos agora definir a convergência das séries numéricas usando a definição da convergência
das sequências numéricas.

%------------------------------------------------------------------------------%
\begin{definicao}{Convergência de Uma Série}{def:convergenia_series}
  Dada uma série numérica\index{Série!convergência}
  \[
    \sum_{k=1}^\infty a_k
  \]
  dizemos que ela \emph{converge} para um número real $L$
  quando sua sequência de somas parciais $(S_n)$ converge para $L$, isso é, \(S_n\to L\).
  Nesse caso podemos escrever
  \[
    \sum_{k=1}^\infty a_k = L
  \]
  e dizemos que $L$ é a \emph{Soma da Série}\index{Soma da série}.
\end{definicao}
%------------------------------------------------------------------------------%

No caso em que a sequência de somas parciais diverge, dizemos que a série \emph{diverge}.

Quando é evidente que os índices da série variam de $1$ até $\infty$,
podemos usar apenas a notação \index{Série!notação}
\[
  \sum a_n   \,.
\]
Entretanto, muitas vezes as séries começam com índices diferentes de 1,
nesses casos os índices precisam ser explicitados, por exemplo
\[
  \sum_{n= 0}^\infty a_n   \qquad \text{ou} \qquad
  \sum_{n=-1}^\infty a_n   \,.
\]

%------------------------------------------------------------------------------%

Para ilustrar os conceitos envolvendo séries infinitas apresentados,
vamos explorar a série
\[
  \sum_{n=1}^{\infty} \frac{1}{2^k}  \,.
\]
Essa série é interessante pois podemos interpretá-la
geometricamente com facilidade.
Note que os primeiros termos da sequência que queremos somar são
\[
  a_1 = \frac{1}{2}   \,,\qquad
  a_2 = \frac{1}{4}   \,,\qquad
  a_3 = \frac{1}{8}   \,,\qquad
  a_4 = \frac{1}{16}  \,,\qquad
  a_5 = \frac{1}{32}  \,.
\]
Podemos imaginar que cada termo é a área de um retângulo e queremos
encontrar a área da figura formada pela união de todos esses retângulos.
A Figura~\ref{fig:soma-1div2n} ilustra geometricamente os cinco primeiros
termos da sequência de somas parciais dessa série e sua soma. O próximo exemplo
efetua os cálculos para mostrar que a série realmente converge e que
sua soma é igual a $1$.

%------------------------------------------------------------------------------%
\begin{figure}
	\centering
  \input{fig_serie_1div2n_sum_01}
	\input{fig_serie_1div2n_sum_02}
	\input{fig_serie_1div2n_sum_03} \\
	\input{fig_serie_1div2n_sum_04}
	% Title: gl2ps_renderer figure
% Creator: GL2PS 1.4.2, (C) 1999-2020 C. Geuzaine
% For: Octave
% CreationDate: Wed Sep 15 16:38:12 2021
\setlength{\unitlength}{1pt}
\begin{picture}(0,0)
\includegraphics[width=5cm,height=5cm]{fig_serie_1div2n_sum_05}
\end{picture}%
\begin{picture}(141,141)(0,0)
\fontsize{18}{0}\selectfont\put(70.8661,41.7874){\makebox(0,0){\textcolor[rgb]{0,0,0}{{$S_{5}$}}}}
\end{picture}

	\input{fig_serie_1div2n_sum}   \\
	\caption{Ilustração das somas parciais da série \(\sum_{n=1}^\infty\frac{1}{2^n}\).}
	\label{fig:soma-1div2n}
\end{figure}
%------------------------------------------------------------------------------%

%------------------------------------------------------------------------------%
\begin{example}
	Calcular a soma os termos da sequência \( a_k = \frac{1}{2^{k}} \).
	\solution
	Queremos calcular a soma da série
	\[
		\sum_{n=1}^{\infty} \frac{1}{2^k}   \,.
	\]
  O gráfico a seguir ilustra a soma dessa série, cada retângulo tem base igual a um
  e altura igual ao termo correspondente da série $a_n$.
  %----------------------------------------------------------------------------%
  \IncludeTexGraphicCentred{fig_serie_1div2n_exemplo_1}[0][-1]
  %----------------------------------------------------------------------------%
	Para analisarmos se a série converge e qual seria o valor da sua soma, vamos construir a
	sequência de somas parciais dessa série:
	\[
		\begin{array}{ll}
			S_1 = a_1                   & = \dfrac{1}{2}                                               \,,\\[3mm]
			S_2 = a_1 + a_2             & = \dfrac{1}{2} + \dfrac{1}{4}                                \,,\\[3mm]
			S_3 = a_1 + a_2 + a_3       & = \dfrac{1}{2} + \dfrac{1}{4} + \dfrac{1}{8}                 \,,\\[3mm]
			S_4 = a_1 + a_2 + a_3 + a_4 & = \dfrac{1}{2} + \dfrac{1}{4} + \dfrac{1}{8} + \dfrac{1}{16} \,.
		\end{array}
	\]
  O próximo gráfico mostra os primeiros termos da sequência de somas parciais.
  %----------------------------------------------------------------------------%
  \IncludeTexGraphicCentred{fig_serie_1div2n_exemplo_S}[-1][-1]
  %----------------------------------------------------------------------------%

  Podemos observar nos gráficos que, como esperado, a série converge para um.
  Porém, não podemos ficar dependentes da evidencia gráfica, que muitas vezes
  é enganosa. Precisamos verificar rigorosamente esse resultado.
  Para isso precisamos identificar a fórmula geral para cada soma parcial.
	Esse é um ponto onde alguns livros escrevem ``\textit{é fácil ver que}''
	para não serem forçados a escrever ``\textit{após muita dor e sofrimento descobrimos que}''
	\[
		S_1 = 1 - \dfrac{1}{2}     \,,\qquad\quad
		S_2 = 1 - \dfrac{1}{4}     \,,\qquad\quad
		S_3 = 1 - \dfrac{1}{8}     \,,\qquad\quad
		S_4 = 1 - \dfrac{1}{16}    \,.
	\]
	Observando o padrão podemos escrever uma conjectura sobre como a fórmula geral deve ser
	\[
		S_n = 1 - \frac{1}{2^{n}} \,.
	\]
	Agora precisamos verificar se ela é verdadeira para todos os valores de $n\in\N$.
	Uma técnica para fazer essa demonstração é utilizar a
	\emph{Indução Finita}\index{Indução!finita}, veja a Seção~\ref{sec:Inducao-Finita}.
	Para isso precisamos primeiro mostrar
	que a fórmula vale para $n=1$. Depois assumimos que ela vale para um $n$ qualquer
	e, partido dessa hipótese, mostrarmos que ela vale para o próximo, $n+1$.

  A demostração que a fórmula vale para $n=1$ pode ser feita por mera avaliação.
  Vamos agora assumir que a fórmula vale para um valor qualquer $n=k-1$
  \begin{equation}
  	\label{eq:hipotese-inducao-exemplo}
  	S_{k-1} = 1 - \frac{1}{2^{k-1}} \,,
  \end{equation}
  queremos mostrar que isso implica que ela também vale para o próximo $n=k$.
  Começamos calculando $S_{k}$
  \[
  	S_{k} = S_{k-1} + a_{k} = S_{k-1} + \frac{1}{2^{k}} \,,
  \]
  substituindo a
  \emph{hipótese de indução}\index{Hipótese de indução}~\eqref{eq:hipotese-inducao-exemplo}
  temos
  \begin{align*}
  	S_{k} & = 1 - \frac{1}{2^{k-1}} + \frac{1}{2^{k}} \\[3mm]
  	      & = 1 - \frac{2}{2^{k}}   + \frac{1}{2^{k}} \\[3mm]
  	      & = 1 - \frac{1}{2^{k}} \,,
  \end{align*}
	que é a fórmula que desejamos demonstrar.

 Tendo uma expressão explícita para os termos da sequência de somas parciais,
 podemos calcular seu limite e determinar a soma da série.
 \begin{align*}
		\lim_{n\to\infty} S_n & = \lim \left(1 - \frac{1}{2^{n}}\right) \\[2mm]
		                      & = \lim 1 -   \lim\frac{1}{2^{n}}        \\[2mm]
		                      & = 1 - 0 = 1    \,.
	\end{align*}
	Portanto, a série converge e podemos escrever
	\[
		\sum_{n=1}^{\infty} \frac{1}{\;2^{n}\;} = 1  \,.
	\]
\end{example}
%------------------------------------------------------------------------------%

\OnlineBox{https://colab.research.google.com/drive/1CE55fvSVp98H4I45LttLroZcLQVljhbZ?usp=sharing}
{
  Código Python para ilustrar algumas séries e suas somas parciais.
}

%------------------------------------------------------------------------------%
%------------------------------------------------------------------------------%
\subsection*{Exemplos de Séries}
%------------------------------------------------------------------------------%
%------------------------------------------------------------------------------%

Apresentamos agora alguns exemplos de séries numéricas importantes.
A série geométrica e a telescópica são exemplos para os quais podemos calcular
a soma da série com relativa facilidade.
Além dessas duas séries, também apresentaremos nessa seção as
séries do inverso do quadrado e inverso do fatorial.

%------------------------------------------------------------------------------%
\begin{definicao}{Série Geométrica}{def:serie-geometrica}
  Sejam $\alpha$ e $r$ números reais com $\alpha\neq 0$, uma série com a forma
  \[
    \sum_{n=1}^\infty \alpha r^{n-1} =
    \alpha + \alpha r + \alpha r^2 + \alpha r^3 + \cdots + \alpha r^{n-1} + \cdots
  \]
  é chamada de \emph{Série Geométrica}\index{Série!geométrica}.
\end{definicao}
%------------------------------------------------------------------------------%

No ensino básico a sequência dos termos dessa série é chamada de
Progressão Geométrica\index{Progressão!geométrica}.
Esse é um exemplo especialmente importante pois conseguimos verificar sua convergência
e calcular sua soma com relativa facilidade.

%------------------------------------------------------------------------------%
\begin{proposicao}{Convergência da Série Geométrica}{pro:covnergencia-serie-geometrica}
  A série geométrica converge se $\abs{r}<1$ e diverge se  $\abs{r}\geq 1$.

  No caso em que a série converge, isso é, se $\abs{r}<1$ a soma da série é dada por
  \begin{equation}
    \label{eq:soma-serie-geometrica}
    \sum_{n=1}^\infty \alpha r^{n-1} = \frac{\alpha}{1-r}  \,.
  \end{equation}
\end{proposicao}
%------------------------------------------------------------------------------%
\begin{proof}
  A sequência de somas parciais da série geométrica é
  \[
    S_n = \sum_{k=1}^n \alpha r^{k-1}  \,.
  \]
  Subtraindo $rS_n$ de $S_n$, obtemos
  \begin{align*}
    S_n(1-r) & = S_n-rS_n                                  \\
             & = \sum_{k=1}^n \alpha r^{k-1}-\sum_{k=1}^n \alpha r^{k}  \\
             & = \left( \alpha
                      + \alpha r
                      + \cdots
                      + \alpha r^{n-1}
                 \right) -
                 \left( \alpha r
                      + \cdots
                      + \alpha r^{n-1}
                      + \alpha r^{n}
                 \right) \\
             & = \alpha - \alpha r^n   \,.
  \end{align*}
  Temos então que
  \[
    S_n = \frac{\alpha \left(1-r^n\right)}{1-r}   \,.
  \]
  Como $r^n\to 0$ se $\abs{r}<1$ e $\abs{r^n}\to\infty$ se $\abs{r}\geq 1$,
  segue que
  \[
    S_n\to \frac{\alpha}{1-r}
  \]
  se $\abs{r}<1$ e $S_n$ diverge se $\abs{r}\geq 1$.
\end{proof}
%------------------------------------------------------------------------------%

Observe que na demonstração todas as manipulações algébricas aconteceram com as
somas parciais, que são somas finitas e não envolvem limites. Só depois de
obtermos uma expressão explícita para $S_n$ é que fizemos $n$ tender para o
infinito e calculamos o limite.

Ao aplicar a proposição acima, cuidado com os índices, por exemplo, a série
\[
  \sum_{n=1}^\infty br^n
\]
também é uma série geométrica, mas a proposição não pode ser aplicada diretamente.
Primeiro precisamos reescrever a série de modo que os índices sejam compatíveis com os da proposição.
Nesse caso,
\[
  \sum_{n=1}^\infty br^n=\sum_{n=1}^\infty brr^{n-1}   \,.
\]
Então podemos tomar $\alpha = br$ e se $\abs{r}<1$ teremos que
\[
  \sum_{n=1}^\infty br^n = \frac{br}{1-r}   \,.
\]

%------------------------------------------------------------------------------%
\begin{example}
  Verifique a convergência da série
  \(
    \sum_{n=1}^\infty \frac{2^{n+1}}{5^n}
  \).

  \solution
  Essa série pode ser reescrita da seguinte maneira
  \[
    \sum_{n=1}^\infty \frac{4}{5}\left(\frac{2}{5}\right)^{n-1}  \,.
  \]
  Então é uma série geométrica com $\alpha =\sfrac{4}{5}$ e $r=\sfrac{2}{5}$.
  Como $\abs{r}=\sfrac{2}{5}<1$, segue do critério de convergência da série geométrica que
  \[
    \sum_{n=1}^\infty \frac{2^{n+1}}{5^n}
      = \frac{\sfrac{4}{5}}{1-\sfrac{2}{5}}
      = \frac{\sfrac{4}{5}}{\sfrac{3}{5}}
      = \frac{4}{3}    \,.
  \]
\end{example}
%------------------------------------------------------------------------------%

%------------------------------------------------------------------------------%
\begin{example}
    Escreva a dízima periódica $0,\overline{d}=0,ddddd\dots$
    (onde $d\neq0$ é um dígito) como a razão de dois números.

    \solution
    \begin{align*}
      0,ddd\,ddd\dots & =\frac{d}{10}+\frac{d}{(10)^2}+\frac{d}{(10)^3}+\cdots           \\
                      & =\frac{d}{10}\left(1+\frac{1}{10}+\frac{1}{(10)^2}+\cdots\right) \\
                      & =\frac{d}{10}\sum_{n=1}^\infty \left(\frac{1}{10}\right)^{n-1}   \\
                      & =\frac{d}{10}\left(\frac{1}{1-\sfrac{1}{10}}\right)              \\
                      & =\frac{d}{10}\left(\frac{1}{0,9}\right)                          \\
                      & =\frac{d}{9}                                                     \,.
    \end{align*}
\end{example}
%------------------------------------------------------------------------------%

Note que o resultado obtido nesse exemplo mostra que $0,999\,999\dots$ é \emph{igual} a 1.

%------------------------------------------------------------------------------%
\begin{definicao}{Série Telescópica}{def:serie_telescopica}
  Dizemos que uma série é \emph{telescópica}\index{Série!telescópica}
  se todos os termos, exceto o primeiro e o último, da $n$-ésima soma parcial se anulam.
\end{definicao}
%------------------------------------------------------------------------------%

Essa definição e seu uso para calcular o limite de uma série são melhor
explicados com a utilização de um exemplo.

%------------------------------------------------------------------------------%
\begin{example}
  Encontre a soma da série
  \(
    \sum_{n=1}^\infty \frac{1}{2n+1} - \frac{1}{2n+3}
  \).

  \solution
  Escrevendo os primeiros temos dessa série percebemos que eles se anulam quando somados
  \begin{align*}
  a_1 & \;=\; \frac{1}{2\cdot 1+1}-\frac{1}{2\cdot 1+3} \;=\; \frac{1}{3}-\frac{1}{5} \,,\\[2mm]
  a_2 & \;=\; \frac{1}{2\cdot 2+1}-\frac{1}{2\cdot 2+3} \;=\; \frac{1}{5}-\frac{1}{7} \,,\\[2mm]
  a_3 & \;=\; \frac{1}{2\cdot 3+1}-\frac{1}{2\cdot 3+3} \;=\; \frac{1}{7}-\frac{1}{9} \,.
  \end{align*}
  Assim, expandindo os termos da série e cancelando a segunda
  fração de um termo com a primeira do termo seguinte temos
  \begin{align*}
  	S_k & = a_1 + a_2 + \cdots + a_k \\
        & = \left( \frac{1}{3}    - \frac{1}{5}    \right)
          + \left( \frac{1}{5}    - \frac{1}{7}    \right)
          + \left( \frac{1}{7}    - \frac{1}{9}    \right) + \cdots
          + \left( \frac{1}{2k+1} - \frac{1}{2k+3} \right) \\[2mm]
  	    & = \frac{1}{3}-\frac{1}{2k+3}  \,.
  \end{align*}
  Dessa forma,
  \[
    \sum_{n=1}^\infty a_n
    = \lim_{k\to \infty} S_k
    = \lim_{k\to \infty} \left( \frac{1}{3}-\frac{1}{2k+3} \right)
    = \frac{1}{3}   \,.
  \]
\end{example}
%------------------------------------------------------------------------------%

%------------------------------------------------------------------------------%
\begin{example}
  Encontre a soma da série
  \(
    \sum_{n=1}^\infty \frac{4}{(4n-3)(4n+1)}
  \).

  \solution
  Primeiro escrevemos o termo geral em somas parciais
  \[
    a_n = \frac{4}{(4n-3)(4n+1)}=\frac{A}{4n-3}+\frac{B}{4n+1}   \,.
  \]
  Então, $A(4n+1)+B(4n-3)=4$. Assim, para $n=1$ e $n=2$, temos o seguinte sistema
  \[
    \systeme[AB]{ \;5A + B = 4, \;9A + 5B = 4}
  \]
  Resolvendo, temos $A=1$ e $B=-1$, ou seja,
  \[
    a_n = \frac{4}{(4n-3)(4n+1)}=\frac{1}{4n-3}-\frac{1}{4n+1}  \,.
  \]
  Escrevendo agora os primeiros temos da série:
  \begin{align*}
    a_1 & \;=\; \frac{1}{4\cdot 1-3}-\frac{1}{4\cdot 1+1} \;=\; \frac{1}{1}-\frac{1}{5}   \,,\\[2mm]
    a_2 & \;=\; \frac{1}{4\cdot 2-3}-\frac{1}{4\cdot 2+1} \;=\; \frac{1}{5}-\frac{1}{9}   \,,\\[2mm]
    a_3 & \;=\; \frac{1}{4\cdot 3-3}-\frac{1}{4\cdot 3+1} \;=\; \frac{1}{9}-\frac{1}{13}  \,.
  \end{align*}
  Assim, expandindo os termos da série e cancelando a segunda fração
  de um termo com a primeira do termo seguinte temos
  \begin{align*}
    S_k & = \left(        1       - \frac{1}{5}    \right)
          + \left( \frac{1}{5}    - \frac{1}{9}    \right) + \cdots
          + \left( \frac{1}{4k-3} - \frac{1}{4k+1} \right) \\[2mm]
        & = 1-\frac{1}{4k+1}  \,.
  \end{align*}
  Dessa forma,
  \[
    \sum_{n=1}^\infty \frac{4}{(4n-3)(4n+1)}
    = \lim_{k\to \infty} S_k
    = 1   \,.
  \]
\end{example}
%------------------------------------------------------------------------------%

Dada uma sequência $(x_n)$ qualquer, podemos construir uma série $\sum a_n$
cujo termo geral é dado por
\[
  a_1 = x_1  \,, \qquad a_{n+1} = x_{n+1}-x_n
\]
para todo $n\in \N$.
Então a sequência de somas parciais é dada por
\begin{align*}
	S_n & = a_1 + a_2 + \cdots + a_n  \\
	    & = {x_1} + ({x_2}-{x_1})
                + ({x_3}-{x_2})
                + \cdots
                + (x_n-{x_{n-1}}) \\
	    & = x_n   \,.
\end{align*}
Assim, se a sequência $(x_n)$ converge para $L$, a série $\sum a_n$ também converge para $L$.
Por outro lado, se a sequência $(x_n)$ diverge, a série também diverge.
Essa é uma forma de construir uma infinidade de séries telescópicas.

%------------------------------------------------------------------------------%
%------------------------------------------------------------------------------%
Apresentamos a seguir as Séries do Inverso do Quadrado e do Inverso do Fatorial
e suas somas.

%------------------------------------------------------------------------------%
\begin{definicao}{Série do Inverso do Quadrado}{def:serie-inverso-quadrado}
	A série
  \[
    \sum_{n=1}^\infty \frac{1}{\,n^2\,} = \frac{\,\pi^2}{6} \approx 1,6449
  \]
  é conhecida como \emph{Série do Inverso do Quadrado}\index{Série!inverso do quadrado}.
\end{definicao}
%------------------------------------------------------------------------------%

%------------------------------------------------------------------------------%
\begin{definicao}{Série do Inverso do Fatorial}{def:serie-inverso-fatorial}
	A série
  \[
   \sum_{n=0}^\infty \frac{1}{\,n!\,} = e = \approx 2,7183
  \]
	é chamada de \emph{Série do Inverso do Fatorial}\index{Série!inverso do fatorial}.
\end{definicao}
%------------------------------------------------------------------------------%

Para calcular as somas dessas séries precisamos de resultados que serão
apresentados mais a frente neste capítulo.
Vamos mostrar que a Série do Inverso do Quadrado Converge no
Exemplo~\ref{ex:convergencia-serie-inverso-quadrado}
e calculamos seu valor na subseção sobre o problema de
Basileia~\ref{subsec:problema-Basileia}.
A verificação da convergência da Série do Inverso do Fatorial
está no Exemplo~\ref{ex:convergencia-serie-inverso-fatorial}
e o cálculo da sua soma é feito no
Exemplo~\ref{ex:soma-serie-inverso-fatorial}.

%------------------------------------------------------------------------------%
